\documentclass[../../main.tex]{subfiles}

\begin{document}

\section{熵准则}\label{section:ch11-熵准则}

在 \S11.2 中研究 Riemann 问题时，我们取 Lax 熵条件
\begin{align}
\lambda_k(\mathbf{u}_r)<\sigma(\mathbf{u}_r,\mathbf{u}_l)<\lambda_k(\mathbf{u}_l)
\label{align:ch11-Lax熵条件1}
\end{align}
对某个 $k\in\{1,\cdots,m\}$ 作为容许激波的选择准则。人们对发现各种其他数学正确且物理适当的熵条件有浓厚兴趣，旨在将其应用于守恒律组更复杂的积分解，以获得唯一性准则、关于容许间断的更多信息等。

一个一般原理（我们已在 \S4.5.1 对标量守恒律和 \S10.1 对 Hamilton--Jacobi 方程见过实例）是：物理和数学正确的解应作为正则化方程组
\begin{align}
\mathbf{u}^\varepsilon_t+\mathbf{F}(\mathbf{u}^\varepsilon)_x-\varepsilon\mathbf{u}^\varepsilon_{xx}=0 \quad \text{在 } \mathbb{R}\times(0,\infty) \text{ 中}
\label{align:ch11-粘性正则化}
\end{align}
的解当 $\varepsilon\to 0$ 时的极限出现。想法是将项``$\varepsilon\mathbf{u}^\varepsilon_{xx}$''解释为提供小粘性效应，它将``抹平''尖锐激波。希望是在 $\varepsilon\to 0$ 的极限中研究问题\eqref{align:ch11-粘性正则化} 的各方面，从而发现更一般的熵准则，以补充 Lax 条件\eqref{align:ch11-Lax熵条件1}。

\subsection{消失粘性，行波}\label{subsection:ch11-消失粘性行波}

我们首先通过寻找具有形式
\begin{align}
\mathbf{u}^\varepsilon(x,t)=\mathbf{v}\left(\frac{x-\sigma t}{\varepsilon}\right) \quad (x\in\mathbb{R},\;t\geqslant 0)
\label{align:ch11-行波形式粘性}
\end{align}
的行波解来开始研究抛物方程组\eqref{align:ch11-粘性正则化}，其中速度 $\sigma$ 和轮廓 $\mathbf{v}$ 待求。将\eqref{align:ch11-行波形式粘性} 代入\eqref{align:ch11-粘性正则化}，我们发现 $\mathbf{v}:\mathbb{R}\to\mathbb{R}^m$，$\mathbf{v}=\mathbf{v}(s)$，必须求解 ODE
\begin{align}
\ddot{\mathbf{v}}=-\sigma\dot{\mathbf{v}}+D\mathbf{F}(\mathbf{v})\dot{\mathbf{v}} \quad (\cdot=\tfrac{\mathrm{d}}{\mathrm{d}s}).
\label{align:ch11-行波ODE}
\end{align}

现在假设 $\mathbf{u}_l,\mathbf{u}_r\in\mathbb{R}^m$ 给定，且
\begin{align}
\lim_{s\to-\infty}\mathbf{v}=\mathbf{u}_l,\quad \lim_{s\to+\infty}\mathbf{v}=\mathbf{u}_r,\quad \lim_{s\to\pm\infty}\dot{\mathbf{v}}=0.
\label{align:ch11-行波渐近}
\end{align}
则由\eqref{align:ch11-行波形式粘性} 推出
\begin{align}
\lim_{\varepsilon\to 0}\mathbf{u}^\varepsilon(x,t)=
\begin{cases}
\mathbf{u}_l & \text{若 } x<\sigma t,\\
\mathbf{u}_r & \text{若 } x>\sigma t.
\end{cases}
\label{align:ch11-行波极限}
\end{align}
因此\eqref{align:ch11-粘性正则化} 的解当 $\varepsilon\to 0$ 时的极限给出连接状态 $\mathbf{u}_l,\mathbf{u}_r$ 的激波。现在的计划是仔细研究 $\sigma$ 和 $\mathbf{v}$ 的形式，从而收集关于由\eqref{align:ch11-行波极限} 确定的激波结构的更详细信息。

首要问题是是否真的存在 $\sigma$ 和 $\mathbf{v}$ 求解\eqref{align:ch11-行波ODE}、\eqref{align:ch11-行波渐近}。积分\eqref{align:ch11-行波ODE}，推出
\begin{align}
\dot{\mathbf{v}}=\mathbf{F}(\mathbf{v})-\sigma\mathbf{v}+\mathbf{c}
\label{align:ch11-行波积分}
\end{align}
对某常数 $\mathbf{c}\in\mathbb{R}^m$。由\eqref{align:ch11-行波渐近} 得出
\begin{align*}
\mathbf{F}(\mathbf{u}_l)-\sigma\mathbf{u}_l+\mathbf{c}=\mathbf{F}(\mathbf{u}_r)-\sigma\mathbf{u}_r+\mathbf{c}.
\end{align*}
因此
\begin{align}
\mathbf{F}(\mathbf{u}_l)-\mathbf{F}(\mathbf{u}_r)=\sigma(\mathbf{u}_l-\mathbf{u}_r).
\label{align:ch11-行波RH}
\end{align}
由\eqref{align:ch11-行波渐近}、\eqref{align:ch11-行波积分} 和\eqref{align:ch11-行波RH}，ODE \eqref{align:ch11-行波积分} 变为
\begin{align}
\dot{\mathbf{v}}=\mathbf{F}(\mathbf{v})-\mathbf{F}(\mathbf{u}_l)-\sigma(\mathbf{v}-\mathbf{u}_l).
\label{align:ch11-行波ODE2}
\end{align}

现在将左状态 $\mathbf{u}_l$ 视为给定，假设我们试图构造连接 $\mathbf{u}_l$ 到附近状态 $\mathbf{u}_r$ 的行波。由\eqref{align:ch11-行波RH}，必然 $\mathbf{u}_r\in S_k(\mathbf{u}_l)$ 对某个 $k\in\{1,\cdots,m\}$ 且 $\sigma=\sigma(\mathbf{u}_r,\mathbf{u}_l)$。

\begin{theorem}[真非线性系统的行波存在性。]\label{theorem:ch11-真非线性系统行波存在性}假设对 $k=1,\cdots,m$，对 $(\lambda_k,\mathbf{r}_k)$ 是真非线性的。设 $\mathbf{u}_r$ 选取得充分接近 $\mathbf{u}_l$。则存在连接 $\mathbf{u}_l$ 到 $\mathbf{u}_r$ 的\eqref{align:ch11-粘性正则化} 的行波解，当且仅当
\begin{align}
\mathbf{u}_r\in S_k^-(\mathbf{u}_l)
\label{align:ch11-行波存在条件}
\end{align}
对某个 $k\in\{1,\cdots,m\}$。
\end{theorem}

\begin{proof}
1. 首先假设 $\sigma$ 和 $\mathbf{v}$ 求解\eqref{align:ch11-行波ODE}、\eqref{align:ch11-行波渐近}。如前所述，必然 $\mathbf{u}_r\in S_k(\mathbf{u}_l)$ 对某个 $k$，$\sigma=\sigma(\mathbf{u}_r,\mathbf{u}_l)$。设
\begin{align*}
\mathbf{G}(z)\coloneqq\mathbf{F}(z)-\mathbf{F}(\mathbf{u}_l)-\sigma(z-\mathbf{u}_l).
\end{align*}
ODE \eqref{align:ch11-行波ODE2} 读作 $\dot{\mathbf{v}}=\mathbf{G}(\mathbf{v})$，且 $\mathbf{G}(\mathbf{u}_l)=\mathbf{G}(\mathbf{u}_r)=0$。计算 $D\mathbf{G}(\mathbf{u}_l)=D\mathbf{F}(\mathbf{u}_l)-\sigma I$；故 $D\mathbf{G}$ 在 $\mathbf{u}_l$ 处的特征值为 $\{\lambda_k(\mathbf{u}_l)-\sigma\}_{k=1}^m$，对应右、左特征向量 $\{\mathbf{r}_k,\mathbf{l}_k\}_{k=1}^m$。

2. 由于 $\mathbf{u}_r\in S_k(\mathbf{u}_l)$ 且 $|\mathbf{u}_r-\mathbf{u}_l|$ 小，由 \S11.2.3\cref{theorem:ch11-激波集结构}(iii)，
\begin{align*}
\sigma=\frac{\lambda_k(\mathbf{u}_r)+\lambda_k(\mathbf{u}_l)}{2}+o(|\mathbf{u}_r-\mathbf{u}_l|).
\end{align*}
因此
\begin{align*}
\lambda_k(\mathbf{u}_l)-\sigma=\frac{\lambda_k(\mathbf{u}_l)-\lambda_k(\mathbf{u}_r)}{2}+o(|\mathbf{u}_r-\mathbf{u}_l|).
\end{align*}
为使 ODE 有轨道连接 $s=-\infty$ 处的 $\mathbf{u}_l$ 到 $s=+\infty$ 处的 $\mathbf{u}_r\in S_k(\mathbf{u}_l)$，必须 $\lambda_k(\mathbf{u}_l)-\sigma>0$，否则轨迹当 $s\to-\infty$ 时不收敛到 $\mathbf{u}_l$。因此若 $|\mathbf{u}_r-\mathbf{u}_l|$ 足够小，$\lambda_k(\mathbf{u}_r)<\lambda_k(\mathbf{u}_l)$，即 $\mathbf{u}_r\in S_k^-(\mathbf{u}_l)$。

3. 条件\eqref{align:ch11-行波存在条件} 的充分性证明从略：见 Majda--Pego (J. Diff. Eq. 56 (1985), 229--262)。
\end{proof}

上述结果利用了真非线性假设，但该断言一般成立，只要我们引入 Lax 熵条件\eqref{align:ch11-Lax熵条件1} 的适当变体。因此现在假设 $\mathbf{u}_r\in S_k(\mathbf{u}_l)$ 对某个 $k$，且
\begin{align}
\begin{cases}
\sigma(z,\mathbf{u}_l)>\sigma(\mathbf{u}_r,\mathbf{u}_l) \text{ 对位于曲线 } S_k(\mathbf{u}_l) \text{ 上 } \mathbf{u}_r \text{ 与 } \mathbf{u}_l \text{ 之间的每个 } z.
\end{cases}
\label{align:ch11-Liu熵条件}
\end{align}
条件\eqref{align:ch11-Liu熵条件} 是\textbf{Liu 熵准则}。（注意若 $(\lambda_k,\mathbf{r}_k)$ 真非线性、$\mathbf{u}_r\in S_k^-(\mathbf{u}_l)$ 且 $\mathbf{u}_r$ 充分接近 $\mathbf{u}_l$，此条件自动成立。）

我们可通过再次寻找方程组\eqref{align:ch11-粘性正则化} 的行波解来激发\eqref{align:ch11-Liu熵条件}。假设 $\mathbf{u}_l$ 给定。则只要 $|\mathbf{u}_r-\mathbf{u}_l|$ 足够小，存在行波解 $\mathbf{u}^\varepsilon(x,t)=\mathbf{v}((x-\varsigma t)/\varepsilon)$（$\mathbf{v}$ 求解\eqref{align:ch11-行波ODE}、\eqref{align:ch11-行波渐近}）当且仅当熵条件\eqref{align:ch11-Liu熵条件} 满足。证明见 Conlon (Adv. in Math. 35 (1980), 1--18)。

\begin{example}[$p$-系统的行波。]\label{example:ch11-p系统行波}再次考虑 $p$-系统
\begin{align}
\begin{cases}
u_{1,t}-u_{2,x}=0 & \text{（相容条件）},\\
u_{2,t}-p(u_1)_x=0 & \text{（Newton 定律）},
\end{cases}
\label{align:ch11-p系统行波}
\end{align}
在通常的严格双曲性条件 $p'>0$ 下。我们研究正则化方程组
\begin{align}
\begin{cases}
u^\varepsilon_{1,t}-u^\varepsilon_{2,x}=0,\\
u^\varepsilon_{2,t}-p(u^\varepsilon_1)_x=\varepsilon u^\varepsilon_{2,xx}
\end{cases}
\label{align:ch11-p系统正则化}
\end{align}
的行波解存在性。注意我们仅在第二个方程中加入粘性项。这在物理上是合理的，因为\eqref{align:ch11-p系统行波} 的第一行仅是数学相容条件。

假设 $\mathbf{u}^\varepsilon=\mathbf{v}((x-\varsigma t)/\varepsilon)$ 是\eqref{align:ch11-p系统正则化} 的行波解，满足
\begin{align*}
\lim_{s\to-\infty}\mathbf{v}=\mathbf{u}_l,\quad \lim_{s\to\infty}\mathbf{v}=\mathbf{u}_r,\quad \lim_{s\to\pm\infty}\dot{\mathbf{v}}=0.
\end{align*}
写 $\mathbf{v}=(v_1,v_2)$，由\eqref{align:ch11-p系统正则化} 计算
\begin{align*}
\begin{cases}
-\sigma\dot{v}_1-\dot{v}_2=0,\\
-\sigma\dot{v}_2-p(v_1)\dot{=}\ddot{v}_2.
\end{cases}
\end{align*}
利用渐近条件积分得
\begin{align*}
\begin{cases}
\sigma v_1+v_2=\sigma v_1^l+v_2^l=\sigma v_1^r+v_2^r,\\
\dot{v}_2=\sigma(v_2^l-v_2)+p(v_1^l)-p(v_1)=\sigma(v_2^r-v_2)+p(v_1^r)-p(v_1),
\end{cases}
\end{align*}
其中 $\mathbf{u}_l=(v_1^l,v_2^l)$，$\mathbf{u}_r=(v_1^r,v_2^r)$。特别地，
\begin{align*}
\begin{cases}
\sigma v_1^l+v_2^l=\sigma v_1^r+v_2^r,\\
\sigma v_2^l+p(v_1^l)=\sigma v_2^r+p(v_1^r).
\end{cases}
\end{align*}
解这些方程求 $\sigma$，得
\begin{align}
\sigma^2=\frac{p(v_1^r)-p(v_1^l)}{v_1^r-v_1^l}.
\label{align:ch11-p系统速度}
\end{align}

假设 $v_1^r>v_1^l$。由\eqref{align:ch11-p系统正则化} 可取 $\sigma>0$。此时 Liu 熵准则读作
\begin{align}
\frac{p(z_1)-p(v_1^l)}{z_1-v_1^l}>\frac{p(v_1^r)-p(v_1^l)}{v_1^r-v_1^l}
\label{align:ch11-p系统Liu条件}
\end{align}
对曲线 $S_k(\mathbf{u}_l)$ 上 $\mathbf{u}_l$ 与 $\mathbf{u}_r$ 之间的所有 $z=(z_1,z_2)$。

我们断言 ODE 组\eqref{align:ch11-p系统正则化} 具有渐近边界条件\eqref{align:ch11-p系统行波} 有解当且仅当熵条件\eqref{align:ch11-p系统Liu条件} 成立。为此，合并\eqref{align:ch11-p系统正则化} 的两个方程消去 $v_2$：
\begin{align*}
\dot{v}_1=\frac{p(v_1)-p(v_1^l)}{\sigma}-\sigma(v_1-v_1^l)\coloneqq g(v_1).
\end{align*}
由\eqref{align:ch11-p系统速度}，$g(v_1^l)=0$ 且 $g(v_1^r)=0$。为使 ODE 有解且 $\lim_{s\to-\infty}v_1=v_1^l$，$\lim_{s\to\infty}v_1=v_1^r$，我们需要 $g(z_1)>0$ 对 $v_1^l<z_1<v_1^r$。但这恰是熵准则\eqref{align:ch11-p系统Liu条件}。若 $v_1^r<v_1^l$ 类似计算。
\end{example}

\subsection{熵/熵通量对}\label{subsection:ch11-熵熵通量对}

Lax 和 Liu 的熵准则都对由激波（或粘性逼近的行波）连接的可能左、右状态提供了限制。然而进一步拓宽熵准则以应用于守恒律更一般的积分解是很有意义的。

一个想法是要求积分解满足某些``熵型''不等式。

\begin{definition}\label{definition:ch11-熵熵通量对}
两个光滑函数 $\Phi,\Psi:\mathbb{R}^m\to\mathbb{R}$ 构成守恒律 $\mathbf{u}_t+\mathbf{F}(\mathbf{u})_x=0$ 的\textbf{熵/熵通量对}，只要
\begin{align}
\Phi \text{ 是凸的}
\label{align:ch11-熵凸}
\end{align}
且
\begin{align}
D\Phi(z)D\mathbf{F}(z)=D\Psi(z) \quad (z\in\mathbb{R}^m).
\label{align:ch11-熵通量条件}
\end{align}
\end{definition}

为激发条件\eqref{align:ch11-熵通量条件}，暂时假设 $\mathbf{u}$ 是 PDE 组 $\mathbf{u}_t+\mathbf{F}(\mathbf{u})_x=0$ 的光滑解。计算
\begin{align}
\Phi(\mathbf{u})_t+\Psi(\mathbf{u})_x=D\Phi(\mathbf{u})\cdot\mathbf{u}_t+D\Psi(\mathbf{u})\cdot\mathbf{u}_x=(-D\Phi(\mathbf{u})D\mathbf{F}(\mathbf{u})+D\Psi(\mathbf{u}))\cdot\mathbf{u}_x=0
\label{align:ch11-熵守恒}
\end{align}
由\eqref{align:ch11-熵通量条件}。此计算说量 $\Phi(\mathbf{u})$ 满足标量守恒律，通量为 $\Psi(\mathbf{u})$。

一般来说，由于激波和其他不规则性，\eqref{align:ch11-一维守恒律} 的积分解不够光滑，不能证明上述计算。新想法是将\eqref{align:ch11-熵守恒} 替换为不等式：
\begin{align}
\Phi(\mathbf{u})_t+\Psi(\mathbf{u})_x\leqslant 0 \quad \text{在 } \mathbb{R}\times(0,\infty) \text{ 中}.
\label{align:ch11-熵不等式}
\end{align}

在应用中，$\Phi(\mathbf{u})$ 有时是物理熵的负值，$\Psi(\mathbf{u})$ 是熵通量。因此不等式\eqref{align:ch11-熵不等式} 断言熵按其通量演化，但也可能经历急剧增加，例如沿激波。

我们从此以后严格理解\eqref{align:ch11-熵不等式} 意为
\begin{align}
\int_0^\infty\int_{-\infty}^\infty \Phi(\mathbf{u})v_t+\Psi(\mathbf{u})v_x\,\mathrm{d}x\mathrm{d}t\geqslant 0
\label{align:ch11-熵不等式弱形式}
\end{align}
对每个 $v\in C_c^\infty(\mathbb{R}\times(0,\infty))$，$v\geqslant 0$。

再次考虑初值问题
\begin{align}
\begin{cases}
\mathbf{u}_t+\mathbf{F}(\mathbf{u})_x=0 & \text{在 } \mathbb{R}\times(0,\infty) \text{ 中},\\
\mathbf{u}=\mathbf{g} & \text{在 } \mathbb{R}\times\{t=0\} \text{ 上}.
\end{cases}
\label{align:ch11-初值问题熵}
\end{align}

\begin{definition}\label{definition:ch11-熵解}
我们称 $\mathbf{u}$ 为\eqref{align:ch11-初值问题熵} 的\textbf{熵解}，只要 $\mathbf{u}$ 是积分解且对每个熵/熵通量对 $(\Phi,\Psi)$ 满足不等式\eqref{align:ch11-熵不等式弱形式}。
\end{definition}

现在尝试为一般初值 $\mathbf{g}$ 构造熵解。如 \S11.4.1 中，我们期望这样的``物理正确''解 $\mathbf{u}$ 是逼近粘性问题
\begin{align}
\begin{cases}
\mathbf{u}^\varepsilon_t+\mathbf{F}(\mathbf{u}^\varepsilon)_x-\varepsilon\mathbf{u}^\varepsilon_{xx}=0 & \text{在 } \mathbb{R}\times(0,\infty) \text{ 中},\\
\mathbf{u}^\varepsilon=\mathbf{g} & \text{在 } \mathbb{R}\times\{t=0\} \text{ 上}
\end{cases}
\label{align:ch11-粘性逼近}
\end{align}
的解 $\mathbf{u}^\varepsilon$ 的极限。假设 $\mathbf{u}^\varepsilon$ 是\eqref{align:ch11-粘性逼近} 的光滑解，当 $|x|\to\infty$ 时足够快地趋于零以证明以下计算。进一步假设 $\{\mathbf{u}^\varepsilon\}_{0<\varepsilon\leqslant 1}$ 在 $L^\infty$ 中一致有界，且
\begin{align}
\mathbf{u}^\varepsilon\to\mathbf{u} \quad \text{a.e. 当 } \varepsilon\to 0
\label{align:ch11-ae收敛}
\end{align}
对某极限函数 $\mathbf{u}$。（实际上验证此 a.e. 收敛极其困难。）

\begin{theorem}[熵与消失粘性。]\label{theorem:ch11-熵与消失粘性}函数 $\mathbf{u}$ 是守恒律\eqref{align:ch11-初值问题熵} 的熵解。
\end{theorem}

\begin{proof}
1. 选取任意光滑熵/熵通量对 $(\Phi,\Psi)$。将\eqref{align:ch11-粘性逼近} 左乘 $D\Phi(\mathbf{u}^\varepsilon)$ 并回忆\eqref{align:ch11-熵通量条件}，计算
\begin{align}
\Phi(\mathbf{u}^\varepsilon)_t+\Psi(\mathbf{u}^\varepsilon)_x=\varepsilon D\Phi(\mathbf{u}^\varepsilon)\mathbf{u}^\varepsilon_{xx}=\varepsilon\Phi(\mathbf{u}^\varepsilon)_{xx}-\varepsilon(D^2\Phi(\mathbf{u}^\varepsilon)\mathbf{u}^\varepsilon_x)\cdot\mathbf{u}^\varepsilon_x.
\label{align:ch11-熵粘性推导}
\end{align}
由于 $\Phi$ 凸，
\begin{align}
(D^2\Phi(\mathbf{u}^\varepsilon)\mathbf{u}^\varepsilon_x)\cdot\mathbf{u}^\varepsilon_x\geqslant 0.
\label{align:ch11-凸性}
\end{align}

2. 将\eqref{align:ch11-熵粘性推导} 乘以 $v\in C_c^\infty(\mathbb{R}\times(0,\infty))$，$v\geqslant 0$。分部积分后发现
\begin{align*}
\int_0^\infty\int_{-\infty}^\infty \Phi(\mathbf{u}^\varepsilon)v_t+\Psi(\mathbf{u}^\varepsilon)v_x\,\mathrm{d}x\mathrm{d}t=\int_0^\infty\int_{-\infty}^\infty \varepsilon(D^2\Phi(\mathbf{u}^\varepsilon)\mathbf{u}^\varepsilon_x)\cdot\mathbf{u}^\varepsilon_x v-\varepsilon\Phi(\mathbf{u}^\varepsilon)v_{xx}\,\mathrm{d}x\mathrm{d}t\geqslant-\int_0^\infty\int_{-\infty}^\infty \varepsilon\Phi(\mathbf{u}^\varepsilon)v_{xx}\,\mathrm{d}x\mathrm{d}t,
\end{align*}
最后一个不等式由\eqref{align:ch11-凸性} 和 $v$ 的非负性成立。

现在令 $\varepsilon\to 0$。回忆\eqref{align:ch11-ae收敛} 和控制收敛定理，得
\begin{align*}
\int_0^\infty\int_{-\infty}^\infty \Phi(\mathbf{u})v_t+\Psi(\mathbf{u})v_x\,\mathrm{d}x\mathrm{d}t\geqslant 0.
\end{align*}
因此 $\mathbf{u}$ 验证熵/熵通量不等式\eqref{align:ch11-熵不等式弱形式}。若 $\Phi$ 和 $\Psi$ 不光滑，经逼近后得到相同结论。

3. 最后固定 $\mathbf{v}\in C_c^\infty(\mathbb{R}\times[0,\infty);\mathbb{R}^m)$，将 PDE \eqref{align:ch11-粘性逼近} 与 $\mathbf{v}$ 点积。分部积分后得
\begin{align*}
\int_0^\infty\int_{-\infty}^\infty \mathbf{u}^\varepsilon\cdot\mathbf{v}_t+\mathbf{F}(\mathbf{u}^\varepsilon)\mathbf{v}_x+\varepsilon\mathbf{u}^\varepsilon\cdot\mathbf{v}_{xx}\,\mathrm{d}x\mathrm{d}t+\left.\int_{-\infty}^\infty\mathbf{g}\cdot\mathbf{v}\,\mathrm{d}x\right|_{t=0}=0.
\end{align*}
令 $\varepsilon\to 0$，推出 $\mathbf{u}$ 是\eqref{align:ch11-初值问题熵} 的积分解。
\end{proof}

\begin{example}\label{example:ch11-标量熵}
在标量守恒律情形（即 $m=1$），对任意凸 $\Phi$ 可找到相应的通量函数 $\Psi$，即
\begin{align*}
\Psi(z)\coloneqq\int_{z_0}^z\Phi'(w)F'(w)\,\mathrm{d}w \quad (z\in\mathbb{R}).
\end{align*}
应用见 \S11.4.3。
\end{example}

\begin{example}\label{example:ch11-p系统熵}
对 $p$-系统，$m=2$。为验证\eqref{align:ch11-熵凸}、\eqref{align:ch11-熵通量条件}，需找 $\Phi,\Psi$，$\Phi$ 凸且
\begin{align*}
(\Phi_{z_1},\Phi_{z_2})\begin{pmatrix}0 & -1\\ -p'(z_1) & 0\end{pmatrix}=(\Psi_{z_1},\Psi_{z_2}).
\end{align*}
一个解为
\begin{align*}
\Phi(z)=\frac{z_2^2}{2}+P(z_1),\quad \Psi(z)=-p(z_1)z_2 \quad (z\in\mathbb{R}^2),
\end{align*}
其中 $P'=p$。注意 $\Phi$ 是凸的，因 $p'>0$。更多例子见习题。
\end{example}

\subsection{单个守恒律的唯一性}\label{subsection:ch11-标量守恒律唯一性}

作为 \S11.4.2 思想的进一步说明，现在再次考虑标量守恒律的初值问题
\begin{align}
\begin{cases}
u_t+F(u)_x=0 & \text{在 } \mathbb{R}\times(0,\infty) \text{ 中},\\
u=g & \text{在 } \mathbb{R}\times\{t=0\} \text{ 上}.
\end{cases}
\label{align:ch11-标量初值}
\end{align}
未知量 $u=u(x,t)$ 为实值，$F:\mathbb{R}\to\mathbb{R}$ 是给定的光滑通量函数。

在 \S3.4 中我们仔细研究了问题\eqref{align:ch11-标量初值}，利用 $F$ 严格凸的主要假设导出 Lax--Oleinik 公式（见 \S3.4.2）。现在放弃 $F$ 凸的假设，设计适当的弱解概念。如前，我们引入熵：

\begin{definition}\label{definition:ch11-标量熵对}
两个光滑函数 $\Phi,\Psi:\mathbb{R}\to\mathbb{R}$ 构成守恒律 $u_t+F(u)_x=0$ 的\textbf{熵/熵通量对}，只要
\begin{align}
\Phi \text{ 是凸的}
\label{align:ch11-标量熵凸}
\end{align}
且
\begin{align}
\Phi'(z)F'(z)=\Psi'(z) \quad (z\in\mathbb{R}).
\label{align:ch11-标量熵通量}
\end{align}
\end{definition}

如上面\cref{example:ch11-标量熵} 所述，对每个凸 $\Phi$ 存在相应的通量 $\Psi$。$u$ 的熵条件读作
\begin{align*}
\Phi(u)_t+\Psi(u)_x\leqslant 0 \quad \text{在 } \mathbb{R}\times(0,\infty) \text{ 上}
\end{align*}
对每个熵/熵通量对 $\Phi,\Psi$。这意为
\begin{align}
\int_0^\infty\int_{-\infty}^\infty \Phi(u)v_t+\Psi(u)v_x\,\mathrm{d}x\mathrm{d}t\geqslant 0
\label{align:ch11-标量熵不等式}
\end{align}
对每个 $v\in C_c^\infty(\mathbb{R}\times(0,\infty))$，$v\geqslant 0$。

\begin{definition}\label{definition:ch11-标量熵解}
我们称 $u\in C([0,\infty),L^1(\mathbb{R}))\cap L^\infty(\mathbb{R}\times(0,\infty))$ 为\eqref{align:ch11-标量初值} 的\textbf{熵解}，只要 $u$ 对每个熵/熵通量对 $(\Phi,\Psi)$ 满足不等式\eqref{align:ch11-标量熵不等式} 且 $u(\cdot,t)\to g$ 在 $L^1$ 中当 $t\to 0$。
\end{definition}

\begin{remark}
(i) 此定义取代了 \S3.4.3 中早期的``熵解''定义。

(ii) 在\eqref{align:ch11-标量熵不等式} 中取 $\Phi(z)=\pm z$，$\Psi(z)=\pm F(z)$，推出
\begin{align*}
\int_0^\infty\int_{-\infty}^\infty uv_t+F(u)v_x\,\mathrm{d}x\mathrm{d}t=0
\end{align*}
对所有 $v\geqslant 0$，从而对所有 $v\in C_c^1(\mathbb{R}\times(0,\infty))$。然后可证明
\begin{align*}
\int_0^\infty\int_{-\infty}^\infty uv_t+F(u)v_x\,\mathrm{d}x\mathrm{d}t+\left.\int_{-\infty}^\infty gv\,\mathrm{d}x\right|_{t=0}=0
\end{align*}
对所有 $v\in C_c^1(\mathbb{R}\times[0,\infty))$，因 $u(\cdot,t)\to g$ 在 $L^1$ 中。因此熵解是积分解。
\end{remark}

我们在 \S11.4.2 中讨论了熵解的构造，现在证明唯一性。

\begin{theorem}[单个守恒律熵解的唯一性。]\label{theorem:ch11-标量熵解唯一性}在相差一个零测集的意义下，\eqref{align:ch11-标量初值} 至多存在一个熵解。
\end{theorem}

如 \S10.2\cref{theorem:ch10-粘性解唯一性} 的证明，基本思想是在问题中``加倍变量''。

\begin{proof}
1. 设 $u$ 是\eqref{align:ch11-标量初值} 的熵解。则
\begin{align}
\int_0^\infty\int_{-\infty}^\infty \Phi(u)v_t+\Psi(u)v_x\,\mathrm{d}x\mathrm{d}t\geqslant 0
\label{align:ch11-熵解不等式}
\end{align}
对所有 $v\in C_c^\infty(\mathbb{R}\times(0,\infty))$，$v\geqslant 0$，其中 $\Phi$ 光滑、凸且 $\Psi(z)=\int_{z_0}^z\Phi'(w)F'(w)\,\mathrm{d}w$ 对任意 $z_0$。固定 $\alpha\in\mathbb{R}$，取
\begin{align*}
\Phi_k(z)\coloneqq\beta_k(z-\alpha) \quad (z\in\mathbb{R}),
\end{align*}
其中对每个 $k=1,\cdots$，函数 $\beta_k:\mathbb{R}\to\mathbb{R}$ 光滑、凸，且
\begin{align*}
\begin{cases}
\beta_k(z)\to|z| & \text{一致},\\
\beta_k'(z)\to\operatorname{sgn}(z) & \text{有界、a.e.}.
\end{cases}
\end{align*}
因此 $\Phi_k(z)\to|z-\alpha|$ 一致地对 $z\in\mathbb{R}$。对应于\eqref{align:ch11-熵解不等式} 的通量为
\begin{align*}
\Psi_k(z)=\int_\alpha^z\beta_k'(w-\alpha)F'(w)\,\mathrm{d}w.
\end{align*}
因此对每个 $z$，
\begin{align*}
\Psi_k(z)\to\int_\alpha^z\operatorname{sgn}(w-\alpha)F'(w)\,\mathrm{d}w=\operatorname{sgn}(z-\alpha)(F(z)-F(\alpha)).
\end{align*}
将 $\Phi_k,\Psi_k$ 代入\eqref{align:ch11-熵解不等式} 并令 $k\to\infty$，推出
\begin{align}
\int_0^\infty\int_{-\infty}^\infty |u-\alpha|v_t+\operatorname{sgn}(u-\alpha)(F(u)-F(\alpha))v_x\,\mathrm{d}x\mathrm{d}t\geqslant 0
\label{align:ch11-Kruzhkov不等式}
\end{align}
对每个 $\alpha\in\mathbb{R}$ 和上述 $v$。

2. 接下来设 $\widetilde{u}$ 是另一个熵解。则
\begin{align*}
\int_0^\infty\int_{-\infty}^\infty |\widetilde{u}-\widetilde{\alpha}|\widetilde{v}_s+\operatorname{sgn}(\widetilde{u}-\widetilde{\alpha})(F(\widetilde{u})-F(\widetilde{\alpha}))\widetilde{v}_y\,\mathrm{d}y\mathrm{d}s\geqslant 0,
\end{align*}
其中 $\widetilde{\alpha}\in\mathbb{R}$，$\widetilde{v}\in C_c^\infty(\mathbb{R}\times(0,\infty))$，$\widetilde{v}\geqslant 0$。

现在令 $w\in C_c^\infty(\mathbb{R}\times\mathbb{R}\times(0,\infty)\times(0,\infty))$，$w\geqslant 0$，$w=w(x,y,t,s)$。固定 $(y,s)\in\mathbb{R}\times(0,\infty)$，在\eqref{align:ch11-Kruzhkov不等式} 中取 $\alpha=\widetilde{u}(y,s)$，$v(x,t)=w(x,y,t,s)$。对 $y,s$ 积分，产生不等式
\begin{align}
\int_0^\infty\int_0^\infty\int_{-\infty}^\infty\int_{-\infty}^\infty |u(x,t)-\widetilde{u}(y,s)|w_t+\operatorname{sgn}(u(x,t)-\widetilde{u}(y,s))(F(u(x,t))-F(\widetilde{u}(y,s)))w_x\,\mathrm{d}x\mathrm{d}y\mathrm{d}t\mathrm{d}s\geqslant 0.
\label{align:ch11-加倍1}
\end{align}
类似地，对每个固定的 $(x,t)\in\mathbb{R}\times(0,\infty)$，在\eqref{align:ch11-Kruzhkov不等式} 中取 $\widetilde{\alpha}=u(x,t)$，$\widetilde{v}(y,s)=w(x,y,t,s)$。对 $x,t$ 积分得
\begin{align}
\int_0^\infty\int_0^\infty\int_{-\infty}^\infty\int_{-\infty}^\infty |\widetilde{u}(y,s)-u(x,t)|w_s+\operatorname{sgn}(\widetilde{u}(y,s)-u(x,t))(F(\widetilde{u}(y,s))-F(u(x,t)))w_y\,\mathrm{d}x\mathrm{d}y\mathrm{d}t\mathrm{d}s\geqslant 0.
\label{align:ch11-加倍2}
\end{align}
将\eqref{align:ch11-加倍1}、\eqref{align:ch11-加倍2} 相加：
\begin{align}
\int_0^\infty\int_0^\infty\int_{-\infty}^\infty\int_{-\infty}^\infty |u(x,t)-\widetilde{u}(y,s)|(w_t+w_s)+\operatorname{sgn}(u(x,t)-\widetilde{u}(y,s))(F(u(x,t))-F(\widetilde{u}(y,s)))(w_x+w_y)\,\mathrm{d}x\mathrm{d}y\mathrm{d}t\mathrm{d}s\geqslant 0.
\label{align:ch11-加倍和}
\end{align}

3. 我们如下设计\eqref{align:ch11-加倍和} 中 $w$ 的巧妙选择。选取 $\eta$ 为标准磨光子（\S C.4，$n=1$），写 $\eta_\varepsilon(x)=\frac{1}{\varepsilon}\eta(x/\varepsilon)$。取
\begin{align*}
w(x,y,t,s)\coloneqq\eta_\varepsilon\left(\frac{x-y}{2}\right)\eta_\varepsilon\left(\frac{t-s}{2}\right)\phi\left(\frac{x+y}{2},\frac{t+s}{2}\right),
\end{align*}
其中 $\phi\in C_c^\infty(\mathbb{R}\times(0,\infty))$，$\phi\geqslant 0$。将此 $w$ 代入\eqref{align:ch11-加倍和}，得
{\footnotesize\begin{align}
\begin{split}
\int_0^\infty\int_0^\infty\int_{-\infty}^\infty\int_{-\infty}^\infty \bigg\{&|u(x,t)-\widetilde{u}(y,s)|\phi_t\left(\frac{x+y}{2},\frac{t+s}{2}\right)\\
&+\operatorname{sgn}(u(x,t)-\widetilde{u}(y,s))(F(u(x,t))-F(\widetilde{u}(y,s)))\\
&\quad\times\phi_x\left(\frac{x+y}{2},\frac{t+s}{2}\right)\bigg\}\\
&\quad\times\eta_\varepsilon\left(\frac{x-y}{2}\right)\eta_\varepsilon\left(\frac{t-s}{2}\right)\,\mathrm{d}x\mathrm{d}y\mathrm{d}t\mathrm{d}s\geqslant 0.
\end{split}
\label{align:ch11-磨光}
\end{align}}
作变量替换
\begin{align*}
\bar{x}=\frac{x+y}{2},\quad \bar{t}=\frac{t+s}{2},\quad \bar{y}=\frac{x-y}{2},\quad \bar{s}=\frac{t-s}{2}.
\end{align*}
则\eqref{align:ch11-磨光} 蕴含
\begin{align*}
\int_{-\infty}^\infty\int_{-\infty}^\infty G(\bar{y},\bar{s})\eta_\varepsilon(\bar{y})\eta_\varepsilon(\bar{s})\,\mathrm{d}\bar{y}\mathrm{d}\bar{s}\geqslant 0,
\end{align*}
其中
\begin{align*}
\begin{aligned}
G(\bar{y},\bar{s})\coloneqq\int_0^\infty\int_{-\infty}^\infty \,&|u(\bar{x}+\bar{y},\bar{t}+\bar{s})-\widetilde{u}(\bar{x}-\bar{y},\bar{t}-\bar{s})|\phi_t(\bar{x},\bar{t})\\
&+\operatorname{sgn}(u(\bar{x}+\bar{y},\bar{t}+\bar{s})-\widetilde{u}(\bar{x}-\bar{y},\bar{t}-\bar{s}))\\
&\quad\times(F(u(\bar{x}+\bar{y},\bar{t}+\bar{s}))-F(\widetilde{u}(\bar{x}-\bar{y},\bar{t}-\bar{s})))\phi_x(\bar{x},\bar{t})\,\mathrm{d}\bar{x}\mathrm{d}\bar{t}.
\end{aligned}
\end{align*}
现在 $u(\bar{x}+\bar{y},\bar{t}+\bar{s})\to u(\bar{x},\bar{t})$，$\widetilde{u}(\bar{x}-\bar{y},\bar{t}-\bar{s})\to\widetilde{u}(\bar{x},\bar{t})$ 在 $L^1_{\mathrm{loc}}$ 中当 $\bar{y},\bar{s}\to 0$。由于映射 $(a,b)\mapsto|a-b|$，$\operatorname{sgn}(a-b)(F(a)-F(b))$ 是 Lipschitz 连续的，令 $\varepsilon\to 0$ 推出
\begin{align*}
\int_0^\infty\int_{-\infty}^\infty |u(\bar{x},\bar{t})-\widetilde{u}(\bar{x},\bar{t})|\phi_t(\bar{x},\bar{t})+\operatorname{sgn}(u(\bar{x},\bar{t})-\widetilde{u}(\bar{x},\bar{t}))(F(u(\bar{x},\bar{t}))-F(\widetilde{u}(\bar{x},\bar{t})))\phi_x(\bar{x},\bar{t})\,\mathrm{d}\bar{x}\mathrm{d}\bar{t}\geqslant 0.
\end{align*}
改写 $x=\bar{x}$，$t=\bar{t}$，因此
\begin{align}
\int_0^\infty\int_{-\infty}^\infty a(x,t)\phi_t(x,t)+b(x,t)\phi_x(x,t)\,\mathrm{d}x\mathrm{d}t\geqslant 0,
\label{align:ch11-最终不等式}
\end{align}
其中
\begin{align*}
\begin{cases}
a(x,t)\coloneqq|u(x,t)-\widetilde{u}(x,t)|,\\
b(x,t)\coloneqq\operatorname{sgn}(u(x,t)-\widetilde{u}(x,t))(F(u(x,t))-F(\widetilde{u}(x,t))).
\end{cases}
\end{align*}

4. 现在利用不等式\eqref{align:ch11-最终不等式} 建立 $L^1$ 压缩不等式
\begin{align}
\int_{-\infty}^\infty |u(x,t)-\widetilde{u}(x,t)|\,\mathrm{d}x\leqslant\int_{-\infty}^\infty |u(x,s)-\widetilde{u}(x,s)|\,\mathrm{d}x
\label{align:ch11-L1压缩}
\end{align}
对 a.e. $0\leqslant s\leqslant t$。

为证明此断言，取 $0<s<t$，$r>0$，在\eqref{align:ch11-最终不等式} 中取 $\phi(x,t)=\alpha(x)\beta(t)$，其中
\begin{align*}
\begin{cases}
\alpha:\mathbb{R}\to\mathbb{R} \text{ 光滑},\\
\alpha(x)=1 \text{ 若 } |x|\leqslant r,\;\alpha(x)=0 \text{ 若 } |x|\geqslant r+1,\\
|\alpha'(x)|\leqslant 2,
\end{cases}
\end{align*}
且
\begin{align*}
\begin{cases}
\beta:\mathbb{R}\to\mathbb{R} \text{ Lipschitz},\\
\beta(\tau)=0 \text{ 若 } 0\leqslant\tau\leqslant s \text{ 或 } \tau\geqslant t+\delta,\\
\beta(\tau)=1 \text{ 若 } s+\delta\leqslant\tau\leqslant t,\\
\beta \text{ 在 } [s,s+\delta] \text{ 和 } [t,t+\delta] \text{ 上线性},
\end{cases}
\end{align*}
其中 $0<\delta<t-s$。推出
\begin{align*}
\frac{1}{\delta}\int_s^{s+\delta}\int_{-\infty}^\infty a(x,\tau)\alpha(x)\,\mathrm{d}x\mathrm{d}\tau\geqslant\frac{1}{\delta}\int_t^{t+\delta}\int_{-\infty}^\infty a(x,\tau)\alpha(x)\,\mathrm{d}x\mathrm{d}\tau-\int_s^{t+\delta}\int_{\{r\leqslant|x|\leqslant r+1\}}b(x,\tau)\alpha'(x)\beta(\tau)\,\mathrm{d}x\mathrm{d}\tau.
\end{align*}
令 $r\to\infty$：
\begin{align*}
\frac{1}{\delta}\int_t^{t+\delta}\int_{-\infty}^\infty a(x,\tau)\,\mathrm{d}x\mathrm{d}\tau\leqslant\frac{1}{\delta}\int_s^{s+\delta}\int_{-\infty}^\infty a(x,\tau)\,\mathrm{d}x\mathrm{d}\tau.
\end{align*}
再令 $\delta\to 0$，推出\eqref{align:ch11-L1压缩} 对 a.e. $0\leqslant s\leqslant t$。

5. 由\eqref{align:ch11-L1压缩} 以及 $u(\cdot,t),\widetilde{u}(\cdot,t)\to g$ 在 $L^1$ 中当 $t\to 0$，我们最终得出 $u=\widetilde{u}$ a.e.。
\end{proof}

\end{document}